\section{Conclusão e trabalhos futuros}

Na base de \rnLinhas{} pedidos de empréstimo, uma rede neural com padronização e parada antecipada acerta \rnAccRede~$\pm$~\rnSdRede\%: supera a regressão logística, um neurônio só, nas 10 seeds, e a regra de um corte em \rnSeedsRedeRegra{} delas. Sem padronização, ela não aprende (\rnSemPadAcc\%). Com a paciência padrão de 10 épocas e taxa de aprendizado de $10^{-3}$, pararia cedo demais (\rnParadaAccLentaDez\% de validação). Mesmo bem ajustada, a rede perde para a árvore de decisão (\rnAccArv\%) em todas as seeds, treina cerca de \rnRazaoTreino{} vezes mais devagar e exige uma busca de hiperparâmetros \rnRazaoBusca{} vezes mais longa. O motivo é a forma da regra que gerou os rótulos, um degrau no score de crédito: a árvore o representa com uma pergunta, e a rede só o aproxima com uma rampa, onde ficam seus erros. Nesta base, a rede funciona, mas não compensa.

Dois caminhos seguem do que foi visto: repetir a comparação em dados de crédito reais, onde a relação com o score é gradual, e separar a variação causada pelos pesos iniciais da rede da variação causada pela divisão dos dados.
